\documentclass[11pt, a4paper]{article}

\usepackage{fancyhdr}
\usepackage{amsmath}
\usepackage{amssymb}
\usepackage{amsthm}
\usepackage{geometry}
\usepackage{enumerate}
\usepackage{algorithm}
\usepackage{algpseudocode}
\usepackage[hidelinks]{hyperref}

\geometry{margin=1in, headheight=14pt}

\input{../../public/macros.tex}

\newcommand{\qtext}[1]{\begingroup\itshape #1\par\endgroup\smallskip}
\newcommand{\OPT}{\operatorname{OPT}}
\newcommand{\DC}{\operatorname{DC}}
\newcommand{\with}[1]{\operatorname{w}_{#1}}  % w_0, w_1, w_2 : the constrained minima
\newcommand{\Tv}{T_{v}}

\theoremstyle{plain}
\newtheorem{theorem}{Theorem}
\newtheorem{lemma}[theorem]{Lemma}
\newtheorem{corollary}[theorem]{Corollary}
\theoremstyle{definition}
\newtheorem{definition}[theorem]{Definition}
\renewcommand{\qedsymbol}{$\blacksquare$}

\begin{document}
\pagestyle{fancy}
\fancyhf{}
\fancyhead[R]{Optimal Troupe Stationing}

\begin{center}
    {\LARGE\bfseries Optimal Military Troupe Stationing on a Tree}\\[4pt]
    {\large Design and Analysis of Algorithms Project Report}\\[10pt]
    {\normalsize Ramdas Singh\ (\texttt{bmat2423}) \qquad Sathvik H S\ (\texttt{bmat2428})}
\end{center}
\vspace{1em}

\section{Problem Statement}

\qtext{Country XYZ has $n$ important towns connected by roads. There are exactly
$n-1$ roads, and every town is reachable from every other town by a unique
sequence of roads. Two towns are called neighbors if a direct road connects them.
The king decides to station military troupes in the towns. Each town can be
assigned $0$, $1$, or $2$ troupes, subject to the condition that every town
assigned $0$ troupes has at least one neighboring town assigned $2$ troupes.
Design an algorithm to assign troupes to towns so that the total number of
troupes assigned is minimized.}

\noindent A connected graph on $n$ vertices with exactly $n-1$ edges, in which
every pair of vertices is joined by a unique path, is precisely a \emph{tree}.
The towns and roads therefore form a tree $T = (V, E)$ with $\abs{V} = n$.

A \emph{troupe assignment}, or simply an \emph{assignment}, is a map
$f : V \to \{0, 1, 2\}$. An assignment $f$ is \emph{feasible} if, for every town
$v$ with $f(v) = 0$, at least one neighbor $u$ of $v$ has $f(u) = 2$.

The task is to compute a feasible assignment $f$ that minimizes the total
$\abs{f} \defeq \sum_{v \in V} f(v)$ and to report the minimum. We say that a $0$-town
is \emph{covered} by any neighbor carrying $2$ troupes. Feasibility therefore
requires every $0$-town to be covered. A feasible assignment always exists
because $f \equiv 1$ has no $0$-towns and costs $n$, so the optimum is at most
$n$.

\section{The Model: A Rooted-Tree Dynamic Program}

The assignment condition makes a town dependent only on itself and its
\emph{neighbors}. This locality can be exploited using bottom-up dynamic
programming on a rooted tree. We root the tree at an arbitrary town $r$ and
process it from the leaves upward. For a vertex $v$, let $\Tv$ be the subtree
rooted at $v$, and let $C(v)$ be the set of children of $v$. The subtrees
$\{T_u : u \in C(v)\}$ are pairwise vertex-disjoint, and
$V(\Tv) = \{v\} \cup \bigcup_{u \in C(v)} V(T_u)$. Crucially, the only edge
joining $T_u$ to the rest of $\Tv$ is the edge $vu$. Hence, once $f(v)$ is fixed,
the assignment on each child subtree can be chosen independently. The assignment
condition at any vertex $x \neq v$ in $\Tv$ depends only on the assignment within
the child subtree containing $x$, with one exception. A child $u$ may be covered
by its parent $v$ when $f(v) = 2$ and $f(u) = 0$. This is the optimal-substructure
property that drives the recurrence.

The subtlety is that whether a vertex is covered may depend on a vertex
\emph{outside} its own subtree, namely its parent. We therefore track enough
information at each vertex to decouple its subtree from the parent's eventual
choice. Specifically, we maintain four quantities per vertex.

\begin{definition}[Internal satisfaction]
Fix a rooted subtree $\Tv$ and an assignment $f$ on $V(\Tv)$. A vertex
$x \in V(\Tv)$ is \emph{internally satisfied} if $f(x) \geq 1$, or if $f(x) = 0$
and some neighbor $y$ of $x$ \emph{lying in $\Tv$} has $f(y) = 2$. Every vertex of
$\Tv$ other than $v$ itself has all its neighbors inside $\Tv$, so internal
satisfaction for such a vertex coincides with the assignment condition.
\end{definition}

\begin{definition}[The four states]\label{def:states}
For a vertex $v$, define the following minima over assignments $f$ on $V(\Tv)$,
with the convention $\min\varnothing = +\infty$:
\begin{itemize}
    \item $\with{1}(v)$: the minimum of $\abs{f}$ subject to $f(v) = 1$ and every
    vertex of $\Tv$ being internally satisfied.
    \item $\with{2}(v)$: the minimum of $\abs{f}$ subject to $f(v) = 2$ and every
    vertex of $\Tv$ being internally satisfied.
    \item $\OPT(v)$: the minimum of $\abs{f}$ subject to every vertex of $\Tv$
    being internally satisfied, with $f(v)$ unrestricted.
    \item $\DC(v)$ (``don't care''): the minimum of $\abs{f}$ subject to every
    vertex of $\Tv$, except possibly $v$, being internally satisfied.
\end{itemize}
\end{definition}

\noindent In words, $\OPT(v)$ is the cost of an optimal self-contained solution
on $\Tv$, while $\DC(v)$ is the cost when $v$ is allowed to be an uncovered
$0$-town because its parent will provide the covering $2$. We also write
$\with{0}(v)$ for the variant of $\OPT(v)$ in which $f(v) = 0$ is forced, so that
$\OPT(v) = \min\{\with{0}(v), \with{1}(v), \with{2}(v)\}$.

Because the root $r$ has no parent, every town's assignment condition must be
satisfied inside $T_r = T$. Therefore, the answer to the problem is exactly
$\OPT(r)$.

\subsection{The recurrences}

Writing $S(v) = \sum_{u \in C(v)} \OPT(u)$ for the total optimum over the child
subtrees, the four quantities satisfy
\begin{align}
    \with{1}(v) &= S(v) + 1, \label{eq:w1}\\
    \with{2}(v) &= \sum_{u \in C(v)} \DC(u) \;+\; 2, \label{eq:w2}\\
    \with{0}(v) &= \min_{w \in C(v)} \Big( \with{2}(w) + \!\!\sum_{\substack{u \in C(v)\\ u \neq w}}\!\! \OPT(u) \Big)
                 \;=\; \min_{w \in C(v)} \big( S(v) - \OPT(w) + \with{2}(w) \big), \label{eq:w0}\\
    \OPT(v)     &= \min\big\{\, \with{1}(v),\ \with{2}(v),\ \with{0}(v) \,\big\}, \label{eq:opt}\\
    \DC(v)      &= \min\big\{\, \OPT(v),\ S(v) \,\big\}. \label{eq:dc}
\end{align}
A leaf has $C(v) = \varnothing$, so $S(v) = 0$. The base cases are
$\with{1}(v) = 1$, $\with{2}(v) = 2$, and $\with{0}(v) = +\infty$, giving
$\OPT(v) = 1$ and $\DC(v) = 0$.

\subsection{The algorithm}

A single post-order traversal evaluates the arrays bottom-up. The solution to our
problem is $\OPT[r]$.

\begin{algorithm}[ht]
\caption{Minimum troupe assignment on a tree}
\begin{algorithmic}
\Function{MinTroupes}{tree $T$, root $r$}
    \For{$v$ in a post-order of $T$ rooted at $r$}
        \State $S \gets \sum_{u \in C(v)} \OPT[u]$
        \State $\with{1}[v] \gets S + 1$
        \State $\with{2}[v] \gets \sum_{u \in C(v)} \DC[u] \;+\; 2$
        \State $\with{0} \gets \min_{w \in C(v)} \big( S - \OPT[w] + \with{2}[w] \big)$ \Comment{$+\infty$ if $C(v)=\varnothing$}
        \State $\OPT[v] \gets \min\{\with{1}[v],\ \with{2}[v],\ \with{0}\}$
        \State $\DC[v] \gets \min\{\OPT[v],\ S\}$
    \EndFor
    \State \Return $\OPT[r]$
\EndFunction
\end{algorithmic}
\end{algorithm}

\noindent The three sums and minima at $v$ can be computed in a single pass over
$C(v)$. To recover an optimal \emph{assignment}, rather than merely its cost,
store at each vertex which of the three cases attained the minimum
in~\eqref{eq:opt}, along with the covering child $w$ for the $0$-case. Also store
which case attained the minimum in~\eqref{eq:dc}. A second, top-down pass then
determines $f(v)$ for every vertex $v$, expanding each child in the mode
($\OPT$, $\DC$, or ``must be $2$'') dictated by the stored choice.

\section{Proof of Correctness}

We prove that the recurrences correctly compute the four states. The correctness
of the algorithm then follows because the answer is $\OPT(r)$. Throughout the
proof, we use the structural fact that once $f(v)$ is fixed, the restrictions
$f_u \defeq f\!\restriction_{T_u}$ to the child subtrees are independent, and the
only cross-subtree adjacency is given by each edge $vu$.

\begin{lemma}\label{lem:w1}
$\with{1}(v) = S(v) + 1$.
\end{lemma}
\begin{proof}
Let $f$ be feasible for $\with{1}(v)$, so that $f(v) = 1$ and every vertex of
$\Tv$ is internally satisfied. Since $f(v) = 1 \neq 2$, the vertex $v$ provides no
cover. Therefore, each child $u$ and every other vertex of $T_u$, whose neighbors
all lie in $T_u$, must be internally satisfied within $T_u$. Thus, each $f_u$ is
feasible for $\OPT(u)$, giving $\abs{f} = 1 + \sum_u \abs{f_u} \ge 1 + S(v)$.
Conversely, taking an $\OPT(u)$-optimal assignment $f_u$ on each child subtree and
setting $f(v) = 1$ yields a feasible assignment of cost exactly $1 + S(v)$. The
vertex $v$ is satisfied because $f(v) = 1$, and every other vertex retains its
in-subtree witness. Hence, equality holds.
\end{proof}

\begin{lemma}\label{lem:w2}
$\with{2}(v) = \sum_{u \in C(v)} \DC(u) + 2$.
\end{lemma}
\begin{proof}
Let $f$ be feasible for $\with{2}(v)$, so that $f(v) = 2$ and every vertex of
$\Tv$ is internally satisfied. For a child $u$ with $f(u) = 0$, the covering $2$
may now be $v$ itself, so $u$ need not be internally satisfied within $T_u$. Every
other vertex of $T_u$ must still be internally satisfied. Thus, each $f_u$ is
feasible for $\DC(u)$, and $\abs{f} \ge 2 + \sum_u \DC(u)$. Conversely, choose a
$\DC(u)$-optimal assignment $f_u$ on each child subtree and set $f(v) = 2$. Then
$v$ is satisfied, any child $u$ that was an uncovered $0$-town under $\DC$ is now
covered by $v$, and all other vertices retain their witnesses. Hence, $f$ is
feasible with cost $2 + \sum_u \DC(u)$, proving equality.
\end{proof}

\begin{lemma}\label{lem:w0}
$\with{0}(v) = \min_{w \in C(v)} \big( S(v) - \OPT(w) + \with{2}(w) \big)$, which
is $+\infty$ when $C(v) = \varnothing$.
\end{lemma}
\begin{proof}
Let $f$ be feasible for $\with{0}(v)$, so that $f(v) = 0$ and every vertex of
$\Tv$ is internally satisfied. Since $f(v) = 0$, the vertex $v$ must be covered
from within $\Tv$. Its only in-subtree neighbors are its children, so some child
$w$ must have $f(w) = 2$. For that child, $f_w$ has $f(w) = 2$ and every vertex of
$T_w$ is internally satisfied because the neighbors of $w$ other than $v$ lie in
$T_w$, and $v$ contributes $0$. Thus, $f_w$ is feasible for $\with{2}(w)$. For
every other child $u \neq w$, the value $f(v) = 0$ provides no cover, so $f_u$
must be feasible for $\OPT(u)$. Hence,
$\abs{f} \ge \with{2}(w) + \sum_{u \neq w}\OPT(u) = S(v) - \OPT(w) + \with{2}(w)$.
Minimizing over the designated child $w$ gives the lower bound.

Conversely, choose any child $w$, an optimal assignment for $\with{2}(w)$ on
$T_w$, and an optimal assignment for $\OPT(u)$ on every other child subtree.
Setting $f(v) = 0$ gives a feasible assignment because $v$ is covered by $w$. This
assignment attains the stated bound, proving the reverse inequality. If $v$ has no
children, it cannot be covered, so no feasible assignment exists and
$\with{0}(v) = +\infty$.
\end{proof}

\begin{lemma}\label{lem:optdc}
$\OPT(v) = \min\{\with{0}(v),\with{1}(v),\with{2}(v)\}$ and
$\DC(v) = \min\{\OPT(v),\, S(v)\}$.
\end{lemma}
\begin{proof}
The first identity is immediate. A feasible assignment $f$ for $\OPT(v)$ has
$f(v) \in \{0, 1, 2\}$, and conditioning on that value gives exactly the three
constrained minima.

For $\DC(v)$, the vertex $v$ is exempt from the internal satisfaction
requirement. If $f(v) \ge 1$, or if $f(v) = 0$ and has an in-subtree covering
child, then every vertex of $\Tv$ is internally satisfied, and the cheapest such
assignment costs $\OPT(v)$. The only remaining possibility is $f(v) = 0$ with no
covering child, in which case $v$ relies on its parent. Since $v$ provides no
cover to its children, each child subtree must be feasible for $\OPT(u)$. The
cheapest such assignment costs $\sum_u \OPT(u) = S(v)$. Taking the minimum of
these two possibilities yields $\DC(v) = \min\{\OPT(v), S(v)\}$.
\end{proof}

\begin{theorem}
\textup{\textsc{MinTroupes}}$(T, r)$ returns the minimum number of troupes over
all feasible assignments of $T$.
\end{theorem}
\begin{proof}
We show by induction on $\abs{V(\Tv)}$ that the computed values
$\OPT[v], \DC[v], \with{1}[v]$, and $\with{2}[v]$ equal the quantities defined in
Definition~\ref{def:states}. \emph{Base case:} A leaf $v$ has
$C(v) = \varnothing$. The recurrences give $\with{1}[v] = 1$, $\with{2}[v] = 2$,
$\with{0} = +\infty$, $\OPT[v] = 1$, and $\DC[v] = 0$. These values match the
definitions directly. A lone $0$-town cannot be covered, a lone $1$-town is the
cheapest internally satisfied assignment, and exempting the leaf allows the
zero-cost assignment $f(v) = 0$. \emph{Inductive step:} For an internal vertex
$v$, every child $u$ has a strictly smaller subtree. By the inductive hypothesis,
$\OPT[u]$ and $\DC[u]$ are correct. Lemmas~\ref{lem:w1}, \ref{lem:w2},
\ref{lem:w0}, and \ref{lem:optdc} then show that the values computed at $v$
from~\eqref{eq:w1}--\eqref{eq:dc} equal the defined quantities. Finally, since the
root $r$ has no parent, internal satisfaction throughout $T_r = T$ is exactly the
feasibility condition for the whole tree. Hence, $\OPT(r)$ is the minimum cost of
a feasible assignment, which is what the algorithm returns.
\end{proof}

\begin{corollary}
The top-down reconstruction outputs a feasible assignment $f$ with
$\abs{f} = \OPT(r)$.
\end{corollary}
\begin{proof}
Each stored choice records a case of the corresponding recurrence that attained
the minimum for a vertex processed in a given mode. Expanding the children in the
modes prescribed by that case, as established in
Lemmas~\ref{lem:w1}--\ref{lem:optdc}, reproduces an assignment achieving the
optimal value at each level. The proofs of those lemmas show that each such
assignment is feasible for its corresponding state. By induction, the assignment
produced at the root is feasible and has cost $\OPT(r)$.
\end{proof}

\section{Complexity Analysis}

Each vertex $v$ is processed once, with work proportional to $1 + \abs{C(v)}$. The
sums $S$ and $\with{2}$, along with the minimum for $\with{0}$, can be computed in
a single scan of the children. Summing over all vertices,
\[
    \sum_{v \in V} \big(1 + \abs{C(v)}\big)
    = n + \sum_{v\in V}\abs{C(v)}
    = n + (n - 1)
    = \mathcal{O}(n),
\]
since every non-root vertex is the child of exactly one parent, so
$\sum_v \abs{C(v)} = \abs{E} = n-1$. Building the rooted structure using BFS or
DFS to determine the parents and post-order is also $\mathcal{O}(n)$. The four
arrays and the explicit stack use $\mathcal{O}(n)$ space. Reconstruction adds
another $\mathcal{O}(n)$ pass. Hence, the algorithm runs in $\mathcal{O}(n)$ time
and uses $\mathcal{O}(n)$ space. This is optimal up to constant factors because
merely reading the tree takes $\Omega(n)$ time.

\section{Implementation}

The algorithm is implemented in Python as the module \texttt{troupes.py}. It
provides the linear-time dynamic program, the top-down reconstruction of an
optimal assignment, a brute-force reference solver, and a verification suite that
cross-checks the two over all labelled trees on up to six towns and over random
trees for larger sizes. The complete source, together with the development
notebook and this report, is available at
\begin{center}
\url{https://github.com/MaxMight01/DAA_Project}.
\end{center}
The accompanying \texttt{README.md} lists the requirements and gives the commands
for running the implementation and its verification suite.

\section{Conclusion}

Modelling the towns as a tree and rooting it reduces the problem to a four-state
dynamic program ($\OPT$, $\DC$, $\with{1}$, $\with{2}$) evaluated in a single
post-order pass. The states precisely encode the one piece of non-local
information required by the assignment condition, namely whether a $0$-town may
rely on its parent for coverage. The resulting algorithm is provably correct by
induction over subtrees and achieves optimal $\mathcal{O}(n)$ time and space
complexity.

\end{document}
